\section{Three connected research questions}
Intellectual Statement II documented two Shanghai field observations: Tencent voluntarily disaggregates its emissions disclosure into checkable categories, and China's head-of-state climate pledge shows sovereign commitments are made publicly and are, in principle, trackable. Disclosure is thus a strategic choice, not a fixed attribute --- yet lower-income countries often lack the resources to disclose as polished as high-capacity states, even when practices are sound. Finance ministries, rating agencies, and investors currently treat disclosure as an externally imposed cost rather than a strategic investment shaped by peer behavior.

This proposal's foundation is Ostrom's finding that groups sustain cooperation through mutual monitoring without a top-down enforcer~\cite{ostrom1990}, treating sovereign climate disclosure as a similarly structured public good untested by current sovereign ESG research.

\textbf{Q1---Economics:} Does higher-quality disclosure causally lower borrowing cost, or does the ESG-rating/cost relationship reward rating existence alone~\cite{crifo2017}?

\textbf{Q2---Computation:} Can an observability-based public goods game show that visibility of peers' contributions alone---with no payoff change---raises voluntary contribution?

\textbf{Q3---Behavior:} Does peer visibility raise genuine transparency, or merely teach actors to perform disclosure (``disclosure theater'')?

These are mutually dependent, yielding one question: \emph{when does peer observability cause genuine, rather than merely performed, transparency among self-interested strategic actors?} Figure~\ref{fig:teaser} traces this dependency across the three disciplines.

\section{Answer to Q1: incentives and intellectual lineage}
\textbf{Provisional answer.} The marginal effect of disclosure \emph{quality} on borrowing cost is negative and distinct from merely holding a rating: countries disclosing more granularly should pay a lower spread than same-rated peers, beyond what Crifo et al.~\cite{crifo2017} attribute to rating level alone---falsifiable if no marginal effect survives controlling for rating.

\textbf{Model.} Players are $N=\{1,\dots,5\}$ countries. Each period $t$, country $i$ chooses $c_{i,t}\in[0,e]$, the share of endowment $e$ invested in disclosure, retaining $e-c_{i,t}$ for domestic priorities. Stage payoff is linear: $\pi_{i,t}=a(e-c_{i,t})+b\sum_{j\in N}c_{j,t}$, with $a>b>a/5$, so private retention dominates individually though full contribution maximizes total payoff. Under low observability a country sees only $\sum_j c_{j,t-1}$; under high observability it sees each peer's individual prior contribution.

\textbf{Solution concept.} The benchmark is Nash equilibrium~\cite{nash1950}. Because $a>b$, contributing nothing dominates regardless of others' actions or observability, so $c_{i,t}^{*}=0$ is the unique prediction---independent of the observability regime, since it changes only information, not payoffs. This leaves open why real actors deviate when observed (Q3).

\textbf{Quantities.} $c_{i,t}^{*}=0$ is \emph{hypothetical}. The deployed game's boundary output---\$432 for zero contribution versus \$372 for constant partial contribution---is \emph{observed}. Crifo et al.'s~\cite{crifo2017} coefficient is \emph{estimated} from external data. The linear structure is \emph{inherited} from standard public-goods theory.

\section{Answer to Q2: method and evaluation}
\textbf{Inputs and algorithm.} Section~2's model is a deterministic six-round problem: fixed peer schedule (sums 10,10,10 low observability; 20,22,24 high), player contribution $c_t\in\{0,\dots,10\}$, $a=4$, $b=2$. Each round the engine computes pool $G_t=c_t+\text{(peer sum)}_t$ and payoff $\pi_t=4(10-c_t)+2G_t$, separating the observability manipulation from the payoff rule.

\textbf{Baseline and metric.} The baseline is the deployed Hugging Face game: one human player, four non-adaptive peers, low observability rounds 1--3, high rounds 4--6. Primary metric: average own contribution $\bar c$ by phase.

\textbf{Outputs and verification.} For $c_t=5$ every round, hand-calculated predictions matched the live game exactly across all six rounds, confirming no arithmetic drift. \emph{Typical} ($c_t=5$ throughout): cumulative payoff \$372, contribution unchanged at 5.00 in both phases---payoff rose only because peers' visible contributions rose. \emph{Boundary} ($c_t=0$ throughout): cumulative payoff \$432, the highest, since contributing nothing dominates regardless of peer behavior---consistent with Section~2's Nash prediction. An invalid non-numeric input was correctly rejected.

\textbf{Smallest feasible experiment.} Comparing these fixed strategies isolates the observability manipulation from behavior: both hold $c_t$ constant, so any outcome change is mechanical---a null-behavior baseline, not yet a test of H1--H3.

\textbf{Adaptive extension (implemented).} The companion Colab notebook adds a strategy where rounds 4--6 respond to the prior round's peer average, swept across sensitivities (0--1.0); at maximum sensitivity contribution rises 0$\to$6 tokens, demonstrating only that the mechanism executes as designed. The notebook also sweeps the boundary-versus-typical comparison across 500 random schedules: boundary strictly dominates typical in 100\% of trials.

\section{Answer to Q3: behavior and competing explanations}
\textbf{Behavioral mechanism.} Following Andreoni and Petrie~\cite{andreoni2004}, the leading account is \emph{image-concern}: removing confidentiality raises giving because contributors care how their choice looks to observers.

\textbf{Competing explanation.} \emph{Conditional cooperation}~\cite{fischbacher2001}: contributors match what they observe others contributing, regardless of whether their own choice is watched.

\textbf{Why the current design cannot discriminate.} Peers are fixed and simulated; the player's own contribution is never shown to anyone. Any high-observability rise is therefore consistent with conditional cooperation by construction and cannot evidence image-concern, since being watched never occurs.

\textbf{Discriminating evidence.} A $2\times2$ design crossing \emph{peer visibility} with \emph{player visibility} (does the player believe their own choice is watched?) would separate the mechanisms: a rise only when watched supports image-concern; a rise whenever peer behavior is visible supports conditional cooperation.

\textbf{Limits and next test.} A full run would only establish the mechanism in a toy game against simulated, no-stakes peers---real sovereign actors may have stronger incentives to fake being watched. Testing with real human participants as peers would give ``player visibility'' an actual audience; equal contribution across all four cells regardless of true visibility would rule out image-concern, pointing policy toward peer-behavior visibility rather than one's own record.

\section{Advanced development and future directions}
This proposal builds on two Nobel-recognized foundations: Ostrom's finding that groups sustain cooperation through mutual monitoring without central enforcement~\cite{ostrom1990,nobelostrom2009}, and Nordhaus's integration of climate change into economic analysis~\cite{nordhaus1992,nobelnordhaus2018}. Together they frame this proposal's question: how self-interested actors sustain collective action on a shared resource without a central enforcer, when the actors are sovereign states. What is new today is cheap, source-grounded AI capable of cross-checking disclosure at scale, making peer observability a feasible policy lever for the first time. Sections~2--4 advance this: disclosure is priced into risk, visibility alone can move contribution, and genuine transparency must be distinguished from performance---together showing observability's policy value depends on which behavioral mechanism drives it.

\begin{table}[h]\caption{From laureate foundations to a new interdisciplinary contribution.}
\begin{tabularx}{\columnwidth}{@{}p{1.6cm}Y@{}}\toprule
\textbf{Horizon}&\textbf{Milestone}\\\midrule
Foundation&Ostrom~\cite{ostrom1990}, Nordhaus~\cite{nordhaus1992}: cooperation, climate as unenforced shared resources.\\
PS1&Observability game; verified against hand calculations (Appendix~A).\\
Next study&Human $2\times2$ design separating image-concern from conditional cooperation.\\
2056&Deployed AI auditing tool adopted by a finance ministry as evidence of genuine transparency.\\\bottomrule
\end{tabularx}\end{table}

\noindent\textbf{2056 statement, Open Science \& SDG Contribution.} By 2056, I aim to be recognized for a computational framework for auditable, observability-based climate-finance institutions, testing whether peer observability substitutes for central enforcement in sustaining genuine sovereign disclosure. \textbf{GitHub/Colab:} \GitHubLink, \ColabLink (commit \texttt{ea48874}). \textbf{Hugging Face:} \DemoLink, supporting \textbf{SDG 4---Quality Education}~\cite{sdg4}.